%\vnegxxs
\section{Related Work}
%\vnegxxs
Operator learning represents an innovative approach to surrogate modeling that maps input functions to output functions. Traditional surrogate modeling methods have typically focused on mapping limited sets of system parameters \eg PDE parameters to output functions \eg PDE solution functions, as demonstrated in numerous research works~\citep{higdon2008computer,zhe2019scalable,li2021scalable,wang2021multi,xing2021residual,xing2021deep,li2022infinite}. Neural operator methods, which leverage neural networks as their foundational architecture, have driven substantial progress in operator learning. The Fourier Neural Operator (FNO)~\citep{li2020fourier} emerged alongside a simpler variant, the Low-rank Neural Operator (LNO)~\citep{kovachki2023neural}, which utilizes a low-rank decomposition of the operator's kernel to enhance computational efficiency. Building on these innovations, researchers introduced the Graph Neural Operator (GNO)~\citep{li2020neural}, which innovatively combines Nystrom approximation with graph neural networks to approximate function convolution. The Multipole Graph Neural Operator (MGNO)~\citep{li2020multipole} introduced a multi-scale kernel decomposition approach, achieving linear computational complexity in convolution calculations. Parallel developments included a multiwavelet-based operator learning model~\citep{gupta2021multiwavelet}, which enhanced precision through fine-grained wavelet representations of the operator's kernel. Another significant contribution emerged with the Deep Operator Net (DeepONet)~\citep{lu2021learning}, which employs a dual-network architecture combining a branch network for input functions and a trunk network for sampling locations. This architecture was later refined into the POD-DeepONet~\citep{lu2022comprehensive}, which improved stability and efficiency by replacing the trunk network with POD (or PCA) bases derived from training data. A survey of neural operators is given in \citep{kovachki2023neural}. Recent advances in operator learning have focused on developing mesh-agnostic and data-assimilation approaches, representing a significant departure from traditional methods~\citep{yin2022continuous, chen2022crom, pilva2022learning, boussif2022magnet, esmaeilzadeh2020meshfreeflownet}. These novel approaches remove the constraint of requiring input and output functions to be sampled on fixed or regular meshes, offering greater flexibility in handling diverse data structures. A key innovation in these methods lies in their reformulation of PDE simulation as an ODE-solving problem. This fundamental shift changes the primary objective from simply predicting solutions based on initial or boundary conditions to capturing the complete dynamic behavior of PDE systems. Training PDE surrogates and neural operators requires extensive high-quality data, which poses a significant challenge due to the resource-intensive nature of physical experiments and sophisticated simulators. To optimize predictive performance while minimizing data collection costs, researchers employ multi-fidelity, multi-resolution modeling~\cite{tang2024multi, tang2023large,li2023infinitefidelity}, and active learning approaches~\citep{pmlr-v151-li22b,li2023multiresolution}.

Our approach aligns with recent multi-physics neural operator research, though it addresses distinct goals and employs different strategies~\citep{mccabe2023multiple, rahman2024pretraining, hao2024dpot}. Unlike existing works focusing on general physics surrogate models using diverse simulation data, we specifically target interactions within individual coupled physical systems. Closely related is the Coupled Multiwavelet Neural Operator (CMWNO)~\citep{xiao2023coupled}, which models coupled PDEs using wavelet-based methods. However, CMWNO has limitations, including rigid decomposition schemes, restricted coupling structures, and difficulty handling multiple interacting processes effectively. Our proposed method builds upon flexible FNO layers, enhancing adaptability and extensibility. The introduced aggregation mechanism integrates seamlessly with any latent-feature neural operator framework and naturally supports complex interactions across multiple physical processes.

%Our approach relates to recent advances in leveraging multi-physics data for neural operator training, though with distinct objectives and methodologies~\citep{mccabe2023multiple, rahman2024pretraining, hao2024dpot}. While recent work has focused on developing physics surrogate foundation models by incorporating data from various physics simulations, our research specifically examines the coupled relationships between interaction processes within individual physical systems. The Coupled Multiwavelet Neural Operator (CMWNO)~\citep{xiao2023coupled} represents the closest parallel to our work, as it similarly aims to model coupled PDEs in complex systems with multiple physical processes. CMWNO employs a specialized approach, decoupling the integral kernel during multiwavelet decomposition and reconstruction procedures in the Wavelet space. However, this method faces several limitations: it requires particular structures and fixed function decomposition schemes, and its simple cross-over model structure presents challenges for systems with more than two processes. Additionally, CMWNO sometimes struggles to capture complex correlations among different processes effectively. Our proposed method builds upon FNO layers while offering greater flexibility and adaptability. A key advantage of our approach lies in its extensibility - it can be integrated with any neural operator learning framework that conducts layer-wise functional transformations. Furthermore, our innovative aggregation method readily accommodates any number of physical processes, making it particularly valuable for modeling complex, multi-process systems.
 



 